\pdfbookmark[2]{On VIII.}{nexus-remarks-viii}
\begin{center}\large On VIII.\end{center}

Here $p$ denotes a positive prime of the form $4n+1$; the notation agrees with that used in the paper \emph{Theoria residuorum biquadraticorum} I, art. 23. Thus
\[
\begin{gathered}
f\equiv1\pdot2\pdot3\dotsfour\tfrac12(p-3)\pdot\tfrac12(p-1)\pmod p\\
p=aa+bb;\quad a\equiv1\pmod4;\quad b\equiv af\pmod p
\end{gathered}
\]
The numbers denoted by $\alpha,\beta$ are determined by the decomposition $p=\alpha\alpha+2\beta\beta$. The column $f$ has been added to the two tables found; some gaps in them have also been filled.

The theorem stated in the text is connected with the biquadratic character of the number $2$. For, since (compare \emph{Theoria resid. biqu.} I, art. 21)
\[
2^{\frac{p-1}4}\equiv f^{\frac12 b}\pmod p
\]
the congruence
\[
b\equiv2m+a-1\pmod8
\]
implies
\SourceLeaf{310}
\[
2^{\frac{p-1}4}\equiv f^{m+\frac{a-1}2}\pmod p
\]
and conversely the former follows from the latter. The latter congruence is easily proved as follows. Let $\mu$ be the number of quadratic residues $\alpha_1$ lying between $0$ and $\frac14p$. Then, by VII,
\[
m=2\mu-\tfrac14(p-1)
\]
and the number of quadratic residues $\alpha_2$ lying between $\frac14p$ and $\frac12p$ is $=\frac14(p-1)-\mu$. If now $p\equiv1\pmod8$, so that the number $2$ is a quadratic residue, the numbers $2\alpha_1$ and $p-2\alpha_2$, taken together, form the same collection as the numbers $\alpha_1$ and $\alpha_2$. Denoting the product of these numbers by $A$, one obtains
\[
2^{\frac{p-1}4}A\equiv(-1)^{\frac14(p-1)-\mu}A\pmod p
\]
and consequently
\[
2^{\frac{p-1}4}\equiv(-1)^\mu\equiv f^{2\mu}\equiv f^{m+\frac14(p-1)}\pmod p
\]
Further, in this case $b\equiv0\pmod4$, and consequently
\[
\frac{p-1}4=(a+1)\frac{a-1}4+\frac{bb}4\equiv2\frac{a-1}4\equiv\frac{a-1}2\pmod4
\]
which gives the congruence to be proved,
\[
2^{\frac{p-1}4}\equiv f^{m+\frac{a-1}2}\pmod p
\]

If, on the other hand, $p\equiv5\pmod8$, so that the number $2$ is a quadratic nonresidue, the numbers $2\alpha_1$ and $p-2\alpha_2$ comprise all the quadratic nonresidues between $0$ and $\frac12p$. Denote their product by $B$, and again denote the product of the numbers $\alpha_1$ and $\alpha_2$ by $A$. Then
\[
f\equiv AB,\quad(-1)^{\frac{p-1}4-\mu}2^{\frac{p-1}4}A\equiv B\pmod p
\]
Squaring these two congruences, taking into account that
\[
ff\equiv-1,\quad2^{\frac{p-1}2}\equiv-1\pmod p
\]
gives
\[
-1\equiv AABB,\quad-AA\equiv BB
\]
and hence $A^4\equiv+1$. Now $A$ is a product of quadratic residues, so $AA$ is a product of biquadratic residues and is therefore itself a biquadratic residue. Thus $AA\equiv+1$ must hold, because $-1$ is a biquadratic nonresidue. It follows that
\SourceLeaf{311}
\[
(-1)^{\frac{p-1}4-\mu}2^{\frac{p-1}4}\equiv AB\equiv f\pmod p
\]
and
\[
2^{\frac{p-1}4}\equiv(-1)^{\mu-1}f\equiv f^{2\mu-1}\equiv f^{m+\frac{p-5}4}\pmod p
\]
Finally, in this case $b\equiv2\pmod4$, and consequently
\[
\frac{p-5}4=(a+1)\frac{a-1}4+\frac{bb-4}4\equiv2\frac{a-1}4\equiv\frac{a-1}2\pmod4
\]
which again gives the congruence to be proved,
\[
2^{\frac{p-1}4}\equiv f^{m+\frac{a-1}2}\pmod p
\]
\BookRule{35mm}
